% -----------------------------------------------------------------------------
% future-main-simple.tex
% スポットインスタンス向けスケジューリング戦略（簡易版）
% 詳細版は future-main.tex を参照
% -----------------------------------------------------------------------------
\documentclass[11pt,a4paper]{ltjsarticle}

\usepackage{luatexja}
\usepackage{amsmath,amssymb}
\usepackage{booktabs}
\usepackage{tabularx}
\usepackage{geometry}

\geometry{margin=27mm}

\newcommand{\Fcomp}{F_{\mathrm{comp}}}
\newcommand{\Pspot}{P_{\mathrm{spot}}}
\newcommand{\E}{\mathbb{E}}

\title{スポットインスタンスの\\スケジューリング戦略（簡易版）}
\author{那須 渚}
\date{2026年10月12日}

\begin{document}
\maketitle

% =============================================================================
\section{やりたいこと}

これまでの研究で，$\Fcomp(d)$（期限 $d$ までインスタンスが中断されない確率）
を測定し，\textbf{夜間が最も中断されにくい}ことがわかった．

しかし，\textbf{測っただけで，それを使って何かを決める仕組みがない．}

そこで次は，$\Fcomp(d)$ を使って

\begin{quote}
\textbf{ジョブをどのリージョンで・どのタイプで・いつ実行するか}
\end{quote}

を自動で決める戦略をつくり，GitHub Actions で効果を確かめる．

% =============================================================================
\section{考え方}

中断されるとジョブは最初からやり直しになる．
つまり\textbf{中断されやすい時間帯を選ぶと，やり直しの分だけ余計にお金がかかる．}

この「結局いくらかかるか」を計算できれば，一番安い条件を選べる．
計算は 3 段階で進む．

% -----------------------------------------------------------------------------
\subsection{(1) 平均何回やれば終わるか}

$p$ を 1 回で成功する確率とする（$= \Fcomp$ の実測値）．
成功するまで繰り返すときの平均回数は，

\begin{equation}
\E[\text{試行回数}] = \frac{1}{p}
\end{equation}

「10 回中 7 回成功するなら，1 回成功させるのに $10/7$ 回必要」という意味である．

\begin{table}[htbp]
\centering
\caption{実測値の例（N.~Virginia / c6a.large / $d = 60$ 分）}
\begin{tabular}{@{}lrr@{}}
\toprule
時間帯 & $p$ & 平均試行回数 $1/p$ \\
\midrule
Night   & $0.98$ & $1.02$ 回 \\
Morning & $0.69$ & $1.45$ 回 \\
\bottomrule
\end{tabular}
\end{table}

% -----------------------------------------------------------------------------
\subsection{(2) 合計で何分ぶん課金されるか}

課金される時間は 2 つに分かれる．

\begin{itemize}
  \item \textbf{成功した 1 回} $\to$ $d$ 分ぶん
  \item \textbf{失敗した回数}（平均 $\frac{1}{p}-1$ 回）$\to$
        中断されるまでの時間 $\bar u$ ぶん
\end{itemize}

$\bar u$ は「中断されたジョブが平均何分走ってから落ちたか」で，
観測データ（\texttt{uptime\_s}）から直接わかる．

\begin{equation}
\E[\text{課金時間}] \;=\; d \;+\; \left(\frac{1}{p}-1\right) \times \bar u
\label{eq:billed}
\end{equation}

\paragraph{例（$d = 60$ 分，$\bar u \approx 12$ 分）.}

\begin{align}
\text{Night}\ (p = 0.98):\quad   & 60 + 0.02 \times 12 = 60.2\ \text{分} \\
\text{Morning}\ (p = 0.69):\quad & 60 + 0.45 \times 12 = 65.4\ \text{分}
\end{align}

Morning は Night より \textbf{5.2 分ぶん余計に}課金される．
これが時間帯を選ぶ意味である．

% -----------------------------------------------------------------------------
\subsection{(3) 金額に直す}

スポット価格 $\Pspot$ は 1 時間あたりなので，$60$ で割って分に合わせる．

\begin{equation}
C = \frac{\Pspot}{60} \times \E[\text{課金時間}]
\end{equation}

$\Pspot = 0.0295$ USD/h のとき，
Night は $0.0296$ USD，Morning は $0.0322$ USD となる．

% =============================================================================
\section{戦略：一番安い条件を選ぶ}

ここまでで「リージョン $r$・タイプ $\tau$・起動時刻 $t$ を決めれば
コストが計算できる」状態になった．あとは\textbf{全部試して一番安いものを選ぶ}．

候補は $4\ \text{リージョン} \times 3\ \text{タイプ} \times 24\ \text{時刻}
= 288\ \text{通り}$ しかないので，全列挙でよい．

ただし，ただ安いだけだと「来週の深夜に実行します」という答えが出てしまう．
そこで次の 2 つの条件をつける．

\begin{table}[htbp]
\centering
\caption{戦略が守る条件}
\begin{tabularx}{\textwidth}{@{}lX@{}}
\toprule
条件 & 意味 \\
\midrule
締切 & $t$ は $t_0$ から $t_0 + (D_j - d_j)$ までの範囲で選ぶ \\
\addlinespace
信頼性 & 成功確率が $p_{\min}$ 未満の条件は選ばない \\
\bottomrule
\end{tabularx}
\end{table}

\paragraph{締切条件（slack）.}
$D_j$ は「いつまでに終わればよいか」，$d_j$ は実行時間なので，
\textbf{遅らせてよい上限}は $D_j - d_j$ になる．これを \textbf{slack} と呼ぶ．

例：$10$ 分のジョブで「$6$ 時間以内でよい」なら，約 $5$ 時間 $50$ 分まで待てる．

slack $= 0$ なら今すぐ実行するしかなく，\textbf{選べるのはリージョンとタイプだけ}
になる．これは既発表の
``region selection can substitute for time-shifting'' に対応する．

\paragraph{信頼性条件.}
期待値はあくまで平均なので，「成功率 $30\%$ でも，何度もやり直せば平均は間に合う」
という計算結果が出てしまうことがある．
しかし 3 回に 2 回失敗する CI は実務で使えない．
そこで $p_{\min}$（例：$0.8$）という下限を直接設ける．

% =============================================================================
\section{何と比べるか}

戦略の効果を示すには，比較相手が要る．次の 5 つを比べる．

\begin{table}[htbp]
\centering
\caption{比較する戦略}
\begin{tabularx}{\textwidth}{@{}llX@{}}
\toprule
戦略 & やること & 役割 \\
\midrule
S0: Naive       & 今すぐ・固定の条件で実行 & GitHub Actions の素の挙動 \\
\addlinespace
S1: Cheapest    & 一番安い価格を選ぶ & 価格だけ見る実務的手法 \\
\addlinespace
S2: Region-only & 今すぐ・リージョンだけ選ぶ & 時間シフトできない場合 \\
\addlinespace
S3: Proposed    & リージョン・タイプ・時刻を全部選ぶ & \textbf{提案手法} \\
\addlinespace
S4: Oracle      & 結果を知った上での最良 & 到達可能な上限 \\
\bottomrule
\end{tabularx}
\end{table}

\paragraph{S2 を入れる理由.}
S2 と S3 を比べると，\textbf{時間をずらした効果}と
\textbf{リージョンを変えた効果}を分けて示せる．
査読でほぼ必ず聞かれる点なので，最初から分けておく．

% =============================================================================
\section{注意：少ないデータを信じすぎない}

$\Fcomp$ は $30$--$42$ 件から推定した値なので，
たまたま高く出たセルがある．
そのまま「一番高いセル」を選ぶと，運が良かっただけの条件を
選んでしまう（winner's curse）．

対策は，推定値そのものではなく
\textbf{$95\%$ 信頼区間の下端}（Wilson の下側信頼限界）を使うこと．

\begin{itemize}
  \item サンプルが少ないほど，この値は大きく下がる
  \item つまり\textbf{データの少ない条件を自動的に避けてくれる}
\end{itemize}

計算は \texttt{statsmodels} の \texttt{proportion\_confint(k, n,
method='wilson')} で一行である．
通常版（S3）と保守版（S3-LCB）の両方を比べれば，それ自体が結果になる．

% =============================================================================
\section{実験のやり方}

\subsection{GitHub Actions の構成}

\begin{table}[htbp]
\centering
\caption{ワークフローの 4 ステップ}
\begin{tabularx}{\textwidth}{@{}llX@{}}
\toprule
ジョブ & 実行場所 & やること \\
\midrule
decide    & GitHub hosted & 戦略に従って $(r, \tau, t)$ を決める \\
\addlinespace
provision & GitHub hosted & 決めた時刻まで待ち，スポットを起動して
                            runner を登録 \\
\addlinespace
workload  & スポット上 & 実際の CI ジョブを走らせる \\
\addlinespace
collect   & GitHub hosted & 結果（成功・中断・コスト・時間）を記録 \\
\bottomrule
\end{tabularx}
\end{table}

中断の検知は，既存の \texttt{start\_instances.py} と同じ方法
（インスタンス内から \texttt{spot/instance-action} を監視）が使える．
検知したら失敗として記録し，provision からやり直す．

\paragraph{注意点.}
runner は必ず \texttt{--ephemeral} で登録すること．
そうしないと中断後に古い runner が残り，次の試行が汚れる．

\subsection{振るパラメータ}

\begin{table}[htbp]
\centering
\caption{実験パラメータ}
\begin{tabularx}{\textwidth}{@{}llX@{}}
\toprule
パラメータ & 水準 & 備考 \\
\midrule
$d_j$（実行時間） & $5, 10, 30, 60$ 分 & 既発表と同じ \\
slack             & $0, 1, 3, 6, 12$ 時間 & $0$ は時間シフト不可 \\
$p_{\min}$        & $0, 0.8, 0.9$ & $0$ は制約なし \\
戦略              & S0--S4, S3-LCB & \\
\bottomrule
\end{tabularx}
\end{table}

\subsection{測るもの}

\begin{itemize}
  \item 実際にかかったコスト（USD）
  \item 完了までの時間
  \item やり直した回数
  \item 締切に間に合った割合
  \item \textbf{式 \eqref{eq:billed} の予測と実測のズレ}
\end{itemize}

最後の項目が重要で，これは\textbf{式が現実を説明できているか}の検証になる．

% =============================================================================
\section{まとめ}

\begin{table}[htbp]
\centering
\caption{使う式は 3 つだけ}
\begin{tabularx}{\textwidth}{@{}lX@{}}
\toprule
式 & 意味 \\
\midrule
$\dfrac{1}{p}$ & 成功するまでの平均回数 \\[8pt]
$d + \left(\dfrac{1}{p}-1\right)\bar u$ & 合計の課金時間（分） \\[8pt]
$\dfrac{\Pspot}{60} \times (\text{上式})$ & 金額に換算 \\
\bottomrule
\end{tabularx}
\end{table}

必要なデータは次の 3 つだけで，\textbf{すべて手元にある．}

\begin{itemize}
  \item $p$ $=$ $\Fcomp$（計算済み）
  \item $\bar u$ $=$ 中断時の \texttt{uptime\_s}（データにある）
  \item $\Pspot$ $=$ スポット価格（実測済み）
\end{itemize}

\paragraph{次にやること.}
まず既存の CSV から上の 3 つの表をつくり，
S0--S4 をシミュレーションで比較する \texttt{strategy.py} を書く．

GitHub Actions の実機実験はその後でよい．
シミュレーションで\textbf{差が出る条件を先に見つけてから}実機に移るほうが，
お金も時間も節約できる．

\vspace{1em}
\noindent
\rule{\textwidth}{0.4pt}

\noindent
\small
詳細版（導出・最適化問題としての定式化・実装コード）は
\texttt{future-main.pdf} を参照．

\end{document}
